%! Addition and Multiplication Rules — Unit 3, Lesson 3.3 — Lesson Plan
% GRADUAL-RELEASE plan, timing 5 / 45 / 5 / 5. One Main Ideas / Notes table: I do / we do in rows
% 1-4 (problems 1-11, ~25 min) -> we do the Guided Practice row (problems 12-15, ~20 min). The
% homework (DeltaMath only) is the individual practice. Traps: 2, 6. Crux: 10. Crux transferred: 15.
\documentclass[10pt]{article}
\usepackage{saar-article}
\usepackage{saar-boxes}
\usepackage{graphicx}   % -article does NOT load graphicx
\graphicspath{{images/}}
\newcommand{\TallMath}[1]{$\displaystyle #1\rule[-1.4em]{0pt}{3.2em}$}
\newcommand{\UnitNumberName}{Unit 3: Categorical Data and Probability \quad}
\newcommand{\LessonNumberName}{Lesson 3.3: Addition and Multiplication Rules}

\begin{document}
\begin{center}
    {\Large\bfseries \CourseName \\
    {\small\bfseries \UnitNumberName \LessonNumberName}}
\end{center}

\begin{tcolorbox}[colback=frost, colframe=cerulean, boxrule=0.9pt, arc=2mm,
    left=2mm, right=2mm, top=1.5mm, bottom=1.5mm]
    \textbf{Primary Objective:} Students will apply the addition rule on a Venn diagram, the
    complement rule for ``at least one,'' and the multiplication rule on a probability tree;
    classify pairs of events as complementary, mutually exclusive, independent, or dependent
    \emph{by computing}; and justify in context why two mutually exclusive events that each have
    some chance are dependent.
    \\[2pt]
    \textbf{Standards:} PS.P.1a, PS.P.1c; AFDA.DA.3e (dependent and independent examples),
    AFDA.DA.3f \quad$\cdot$\quad carries forward Lesson~3.2 (complements, Venn diagrams,
    mutually exclusive named)
    \\[2pt]
    \textbf{Lesson model:} \emph{Gradual release --- I do, we do, you do.} Four instruction rows
    (problems 1--11: you read the definition, mark up the display, and work 1, 3, 5, 7, 9; the
    class works the rest), then the Guided Practice row (12--15) with students holding the pen.
    \textbf{The homework is the individual practice}, a DeltaMath assignment launched at Close
    \& Assign and done outside class. The debrief is a \emph{spoken} cold check and one board item.
    \\[2pt]
    \textbf{Scope:} independence here is the plain definition (one event does not change the
    other's probability) and $P(A \text{ and } B) = P(A)P(B)$. The conditional-probability test
    and its notation are Lesson~3.4's.
\end{tcolorbox}

\renewcommand{\arraystretch}{1.4}
\boxguard
\begin{skillbox}[Learning Targets \& Key Understandings]{goldbox}
\begin{tabularx}{\linewidth}{@{}X|X@{}}
\textbf{Learning targets (student language)} & \textbf{Key understandings (the ``why'')} \\[2pt]
\hline
\vspace{-6pt}
\begin{itemize}[leftmargin=1.1em, nosep, after=\vspace{2pt}]
  \item I can use the \textbf{addition rule} on a Venn diagram, subtracting the overlap.
        \hfill \textit{PS.P.1c}
  \item I can add without subtracting when events are mutually exclusive, and find ``at least
        one'' with the complement rule. \hfill \textit{PS.P.1c}
  \item I can use the \textbf{multiplication rule} on a tree for \textbf{independent events}
        (with replacement) and \textbf{dependent events} (without replacement).
        \hfill \textit{PS.P.1c, AFDA.DA.3e}
  \item I can classify two events as complementary, mutually exclusive, independent, or
        dependent by computing. \hfill \textit{PS.P.1a, AFDA.DA.3f}
\end{itemize}
&
\vspace{-6pt}
\begin{itemize}[leftmargin=1.1em, nosep, after=\vspace{2pt}]
  \item \textbf{Target misconception:} that mutually exclusive means independent --- ``they
        can't happen together, so they have nothing to do with each other.'' \textbf{The rule
        that replaces it:} when both events can happen, the two ideas are opposites. If two
        events with nonzero probability are mutually exclusive, knowing one happened tells you
        the other did not, so they are dependent.
  \item \textbf{Second misconception:} adding for ``or'' and ``at least one'' without removing
        the overlap (problems 2 and~6).
  \item A pair is classified by a computation, not by how the words sound: an overlap of $0$
        against a product that is not $0$ is the whole argument.
  \item Replacement is what decides a tree: the second-stage numbers stay the same when the
        ticket goes back and change when it stays out.
\end{itemize}
\end{tabularx}
\end{skillbox}

\boxguard[24]
\begin{skillbox}[{Vocabulary, Concepts, \& Theorems --- taught directly in the guided notes}]{greenbox}
\small
\renewcommand{\arraystretch}{1.35}
\begin{tabularx}{\linewidth}{@{} >{\bfseries}p{3.4cm} X @{}}
Addition rule & $P(A \text{ or } B) = P(A) + P(B) - P(A \text{ and } B)$. \\
Independent events & One event happening does not change the probability of the other;
       $P(A \text{ and } B) = P(A) \cdot P(B)$. \\
Dependent events & One event happening changes the probability of the other. \\
Multiplication rule & $P(A \text{ and } B) = P(A) \cdot P(B)$ for independent events; for
       dependent events the second factor is $B$'s probability after $A$ happened. \\
Mutually exclusive & Cannot happen together; $P(A \text{ and } B) = 0$ (named in Lesson~3.2,
       printed in row~2). \\
Complementary & Mutually exclusive and $P(A) + P(B) = 1$: exactly one of the two happens. \\
\end{tabularx}
\end{skillbox}

\begin{fixedskillbox}[Lesson at a Glance --- \MeetingLength]{frost}
\renewcommand{\arraystretch}{1.25}
\begin{tabularx}{\linewidth}{@{}l c X X@{}}
\textbf{Phase} & \textbf{Min} & \textbf{Students} & \textbf{Teacher} \\
\hline
Warm-Up & 5 & Decimals, fractions, and the raffle Venn --- silent & Debrief item~3's last
  sentence aloud; leave $24 + 20 = 44$ on the board \\
Guided Notes \& Practice & 45 & Fill the vocabulary box as each term is named; work rows 1--4
  (problems 1--11); then the Guided Practice row, \emph{Rain and the Bus} (12--15), together &
  \textbf{I do / we do, row by row}: read, mark up the display, work 1, 3, 5, 7, 9; the class
  works the rest (25) $\to$ \textbf{we do} Guided Practice 12--15 (20) \\
Debrief & 5 & Pens down --- answer aloud, nothing to fill in & Put problem~15 back on the
  board; ask the cold check \\
Close \& Assign & 5 & Write down the DeltaMath assignment and the due date &
  Announce the DeltaMath assignment, post it with its due date, preview Lesson~3.4 \\
\end{tabularx}
\end{fixedskillbox}

\begin{fixedskillbox}[Warm-Up --- Activate Prior Knowledge (5 min)]{frost}
\begin{tabularx}{\linewidth}{@{}X X@{}}
\begin{minipage}[t]{\linewidth}\vspace{0pt}
    \textbf{The three items and what each seeds}
    \begin{itemize}[leftmargin=1.1em, itemsep=1pt, topsep=2pt]
      \item \textbf{1.} $0.6 + 0.5 - 0.3 = 0.8$, $0.6 \cdot 0.5 = 0.3$, $1 - 0.2 = 0.8$.
            \textbf{Seeds:} problems 1, 4, and~9 --- the same three computations.
      \item \textbf{2.} $\frac{3}{10} \cdot \frac{3}{10} = \frac{9}{100}$ and
            $\frac{3}{10} \cdot \frac{2}{9} = \frac{1}{15}$. \textbf{Seeds:} problems 5 and~7,
            the two trees.
      \item \textbf{3.} The drum's Venn (Lesson~3.2): $24$, $20$, $12$, $P(\text{neither}) =
            0.2$. \textbf{Seeds:} row~1's display and the trap at problem~2.
    \end{itemize}
\end{minipage}
&
\begin{minipage}[t]{\linewidth}\vspace{0pt}
    \textbf{Running it}
    \begin{itemize}[leftmargin=1.1em, itemsep=1pt, topsep=2pt]
      \item \textbf{Debrief item~3's sentence aloud}: $24 + 20 = 44$ tickets in a drum of $40$.
            Leave it on the board; row~1 answers it.
      \item The sentence is the tell. A student who cannot say why the count passes $40$ will
            write Jamal's $1.1$ at problem~2.
      \item Items 1--2 are arithmetic. Do not let them expand.
    \end{itemize}
\end{minipage}
\end{tabularx}
\end{fixedskillbox}

\boxguard
\begin{skillbox}[Hook --- before anyone sits down]{frost}
The hook is on \textbf{page 1 of the notes}, under the vocabulary box, and on the board: one
winning ticket; $N$, sold to a ninth grader; $T$, sold to a twelfth grader; one ticket cannot be
both. \textit{Devon says: ``They can't happen together, so they have nothing to do with each
other. They must be independent.''} Students circle \emph{independent}, \emph{dependent}, or
\emph{can't tell} and write a reason \emph{before} anyone speaks; take the hand vote and write
the split on the board. \textbf{Do not settle it.} Say only that the class votes again at
problem~10, in the \emph{Four Words} row.
\end{skillbox}

\boxguard
\begin{skillbox}[Guided Notes \& Practice --- I do, then we do (45 min)]{frost}
\textbf{This is the whole lesson.} One Main Ideas / Notes table: four instruction rows (problems
1--11), then the Guided Practice row (12--15). In every row: read the definition aloud, mark up
the display, work the first problem of the grid yourself, then hand the rest to the room.
\begin{multicols}{2}
\textbf{I do / we do, row by row (about 25 min).}

\emph{Row 1, Addition Rule (1--2).} Name the \emph{addition rule}. Students box the count of
times an overlap ticket is added ($2$). Work problem~1 ($.6 + .5 - .3 = .8$; $32$ of $40$).
\textbf{Problem~2 is the first trap}: Jamal's $1.1$ cannot be a probability. Five minutes.

\emph{Row 2, Mutually Exclusive (3--4).} Read the printed definition. Work problem~3 ($.2 + .25
= .45$; nothing to subtract). The class does~4 on row~1's Venn: $1 - .2 = .8$, the same answer
two ways. Five minutes.

\emph{Row 3, Trees (5--8).} Name \emph{independent events}, \emph{dependent events}, and the
\emph{multiplication rule}. Work 5 ($\frac{9}{100}$) and 7, filling the four boxes
($\frac{2}{9}, \frac{7}{9}, \frac{3}{9}, \frac{6}{9}$; $\frac{1}{15}$). The class does 6 and~8.
\textbf{Problem~6 is the second trap}: Ana's $\frac{6}{10}$ counts (gold, gold) twice; the answer
is $\frac{51}{100}$. Problem~8 is the AFDA.DA.3e example pair. Seven minutes.

\emph{Row 4, Four Words (9--11) --- the crux.} Read the four tests. Work problem~9: $.3 = (.6)(.5)$,
independent, with an overlap. \textbf{Problem~10 is the crux: re-take the hook vote before anyone
computes.} Then $P(N \text{ and } T) = 0$ against $(.2)(.25) = .05$: dependent, and $T$ drops from
$.25$ to $0$ once $N$ is known. Read the boxed sentence word by word. Problem~11 applies it to a
complementary pair ($.24 \neq 0$). Eight minutes.

\columnbreak
\textbf{We do (about 20 min).} \emph{Rain and the Bus} (12--15), a fresh context. \textbf{Students
hold the pen; you hold the questions.} 12: $.45$, said in context. 13: $.15$ against $.08$,
dependent --- rain makes a late bus more likely. 14: $.04$ and $1 - .64 = .36$. \textbf{Problem~15
is the crux transferred}: take a hand count on Theo first; then $0$ against $.02$, dependent.

Circulate and demand the reason, not the word:
\begin{itemize}[leftmargin=1.1em, nosep]
  \item \textbf{12:} ``Which mornings does that $.45$ count?'' (Rain, late, or both.)
  \item \textbf{13:} ``You wrote two numbers. Are they equal, and what does that tell Leah?''
  \item \textbf{14:} ``What is the opposite of late at least once?'' (On time both mornings.)
  \item \textbf{15:} ``If the bus is late today, what is the chance it is early today?''
\end{itemize}
Pick one paper for the debrief: a problem~15 that shows both numbers and says what a late bus
does to $Y$. \textbf{This circulation is your formative read.}
\end{multicols}
\end{skillbox}

\boxguard[24]
\begin{skillbox}[Debrief --- whole class, spoken (5 min)]{frost}
Pens down, papers up. \textbf{This debrief is spoken.}
\begin{multicols}{2}
\begin{enumerate}[leftmargin=1.3em, nosep, itemsep=3pt]
  \item \textbf{The one board item: problem~15.} Put up the paper you picked ($0$ against $.02$)
        and make a student say it: \emph{they can't happen together, so knowing one tells you
        the other did not --- dependent.}
\end{enumerate}
\columnbreak
\textbf{Check for understanding --- ask it cold.} \emph{``Flip one coin once. Event $H$: heads.
Event $T$: tails. Independent or dependent? Give me the two numbers.''} Correct answer:
\textbf{dependent} --- $P(H \text{ and } T) = 0$, but $P(H)P(T) = .25$. \emph{This replaces the
exit ticket and is the lesson's only crux check: the DeltaMath homework does not ask it.}

\textbf{Formative read.} \textbf{(a)} said dependent with both numbers; \textbf{(b)} said
dependent with no numbers; \textbf{(c)} said independent. \textbf{If pile (c) runs past a third
of the room, open Lesson~3.4 by re-running problem~10.}
\end{multicols}
\end{skillbox}

\boxguard
\begin{skillbox}[Homework (DeltaMath) --- scored 1 to 5, assigned today, due after two study halls]{goldbox}
\textbf{DeltaMath only.} The whole assignment is the DeltaMath set \emph{Lesson 3.3: Addition
and Multiplication Rules}, with these topics: Symbolic Probability Rules $\cdot$ Probability with
Or / And $\cdot$ Compound Probability $\cdot$ Determine Independence from a Table $\cdot$
Probability of Independent \& Dependent Events. Nothing is on paper. \textbf{Scored 1 to 5 from
the DeltaMath percent} (5 $\times$ percent, rounded half up, never below 1 for an assignment
completed), due at the start of the first class after two study halls. It is created
\emph{Hidden} in the \textbf{SAAR} folder for \textbf{SAAR - A}: post it and set its due date at
Close \& Assign, and say the scoring aloud.

\textbf{Preview of the next class.} Lesson~3.4, Conditional Probability, Independence, and
Simulation.

\textbf{Connections \& Big Ideas.} \textit{Unit 3 core idea:} probability rules are counting
done carefully, and whether one event bears on another is decided by computing.
\textit{Standards covered:} PS.P.1a (classify events), PS.P.1c (addition, multiplication, and
complement rules), AFDA.DA.3e (dependent and independent examples), AFDA.DA.3f (classify events
in a problem setting). \textit{Reaches back to:} Lesson~3.2's complements and Venn diagrams.
\textit{Spirals forward to:} Lesson~3.4's conditional-probability test and Lesson~3.5's counting.
\end{skillbox}

\boxguard
\begin{skillbox}[Watch For (while circulating)]{redbox}
\begin{itemize}[leftmargin=1.2em, nosep]
  \item \textbf{``Can't happen together, so independent''} (the hook, problems 10 and~15, the
        cold check --- the target misconception). Probe: \emph{``If you know this one happened,
        what is the chance of the other now?''}
  \item \textbf{Adding without the overlap} (problems 2 and~6; DeltaMath's Probability with Or /
        And). Probe: \emph{``Where do the outcomes in both circles get counted?''}
  \item \textbf{Same second-stage numbers on both trees} (problem~7; DeltaMath's Probability of
        Independent \& Dependent Events). Probe: \emph{``How many tickets are left in the
        bowl?''}
  \item \textbf{No overlap read as ``no relationship''} (problem~11; DeltaMath's Determine
        Independence from a Table). Probe: \emph{``Is $0$ the same number as the product?''}
  \item \textbf{A number with no meaning} (problem~12). Probe: \emph{``Fifty-five percent
        of what?''}
\end{itemize}
Cold-call prompts: \emph{``Name two events at this school that cannot happen together. Are they
independent?''} \quad \emph{``What would make the second draw independent?''}
\end{skillbox}

\boxguard
\begin{skillbox}[Close \& Assign (5 min)]{goldbox}
\textbf{The homework is assigned here and done outside class.} Post the DeltaMath assignment
\emph{Lesson 3.3: Addition and Multiplication Rules} and set its due date: the start of the first
class after two study halls; scored 1 to 5 from the DeltaMath percent. Say what it practises:
the ``or'' and ``and'' rules, compound events, and testing independence by multiplying and
comparing. \textbf{Tell students to start with Determine Independence from a Table if they are
short on time.} Close on what changed: \emph{``You walked in thinking `can't happen
together' meant `nothing to do with each other.' You walk out knowing it means the opposite.''}
\textbf{Preview:} Lesson~3.4 gives the second-branch numbers a name and a notation.
\end{skillbox}

\begin{teachernote}[Warm-Up]
Five minutes, silent, timed. Items 1--2 are the arithmetic of problems 1, 4, 5, 7, and~9, so row
work can move. \textbf{Read item~3's last sentence while collecting}: a student who says the
$12$ tickets were counted twice has already made row~1's argument. Note two who cannot say it
and cold-call them at problem~2.
\end{teachernote}

\begin{teachernote}[Guided Notes \& Practice]
\textbf{Pace it 25 / 20, then 5 for the debrief and 5 to assign the homework.} You work 1, 3, 5,
7, and~9; the class works the rest. The traps are 2 and~6. Rows 1 and~2 must move.
\textbf{Spend the time in row~4.} Re-take the hook vote at problem~10 \emph{before} anyone
computes: students who have committed to \emph{independent} and then see $0$ against $.05$ have
a question the boxed sentence answers. Problem~9 comes first on purpose, so students see an
overlap that \emph{is} independent. Keep the claim exact: the rule needs both events to have some
chance. \textbf{Stay out of $P(B \mid A)$ notation}; say ``the chance of $T$ now,'' and leave the
symbol for Lesson~3.4. If time is short in Guided Practice, say problem~12 aloud and protect 13
and~15.
\end{teachernote}

\begin{teachernote}[Homework]
\textbf{DeltaMath only: \emph{Lesson 3.3: Addition and Multiplication Rules}. Scored 1 to 5 from
the DeltaMath percent, due at the start of the first class after two study halls; assigned at the
close, not started in class.} What each topic practises: \textbf{Symbolic Probability Rules} ---
the addition and complement rules in symbols; \textbf{Probability with Or / And} --- ``or'' with
the overlap subtracted, and ``and''; \textbf{Compound Probability} --- multiplying along two
stages; \textbf{Determine Independence from a Table} --- the multiply-and-compare test,
$P(A \text{ and } B)$ against $P(A)P(B)$; \textbf{Probability of Independent \& Dependent
Events} --- with and without replacement, the second factor changing or not.

\textbf{Reading the results.} Misses on Determine Independence from a Table and Probability of
Independent \& Dependent Events signal the target confusion: a student classifying by how the
words sound, not by the computation. Misses on Probability with Or / And point to the dropped
overlap (problem~2's trap). \textbf{DeltaMath has no item asking whether mutually exclusive events
are independent}, so the crux check is the debrief's spoken cold check (the coin, $0$ against
$.25$) and is reopened at Lesson~3.4's warm-up debrief. If the cold check's pile (c) or the
independence-topic misses run past a third of the room, open Lesson~3.4 with problem~10.
\end{teachernote}
\end{document}
